\documentclass[reqno]{amsart}
\usepackage[T1]{fontenc}
\usepackage{amsmath,amssymb,amsthm}
\usepackage{xr-hyper}
\usepackage[hidelinks]{hyperref}
\externaldocument[di-]{20261125T020100--divisibility-at-inert-primes-in-general__workbook}
\externaldocument[el-]{20261124T223000--elliptic-curves-for-the-proof__workbook}
\externaldocument[fk-]{20261125T009500--fock-model-and-theta-intertwiners__workbook}
\externaldocument[cm-]{20261124T223100--complex-multiplication-and-inert-reduction__workbook}
\externaldocument[fo-]{20261126T020120--frames-at-odd-positions-and-unit-scalars__workbook}

\newcounter{chap}
% Chapter numbers, generated from 20261125T012000--grh-workbook__guide.org
% by scripts/make-chapters.py.  Do not edit: change the order of the
% index and rerun the script (./preview-site.sh and CI do this).
%
% Usage in a sheet, after \newcounter{chap}:
%   % Chapter numbers, generated from 20261125T012000--grh-workbook__guide.org
% by scripts/make-chapters.py.  Do not edit: change the order of the
% index and rerun the script (./preview-site.sh and CI do this).
%
% Usage in a sheet, after \newcounter{chap}:
%   % Chapter numbers, generated from 20261125T012000--grh-workbook__guide.org
% by scripts/make-chapters.py.  Do not edit: change the order of the
% index and rerun the script (./preview-site.sh and CI do this).
%
% Usage in a sheet, after \newcounter{chap}:
%   \input{chapters}\chapnumber{<sheet ID>}
\makeatletter
\providecommand\chapters@set[2]{\@namedef{chapters@#1}{#2}}
\chapters@set{20261125T001000}{1}% zeros-to-ramp-sums
\chapters@set{20261125T001100}{2}% inert-primes-and-estimate-b
\chapters@set{20261125T005000}{3}% galois-fourier-and-averaging
\chapters@set{20261125T005500}{4}% columns-from-a-sphere
\chapters@set{20261125T000000}{5}% slope-method
\chapters@set{20261126T030100}{6}% a-twist-for-each-prime
\chapters@set{20261125T003000}{7}% the-chain-and-the-square-root
\chapters@set{20261125T010000}{8}% rank-lemma-and-zero-estimate
\chapters@set{20261125T010500}{9}% poincare-series-and-near-orthogonality
\chapters@set{20261124T223000}{10}% elliptic-curves-for-the-proof
\chapters@set{20261125T009500}{11}% fock-model-and-theta-intertwiners
\chapters@set{20261124T223100}{12}% complex-multiplication-and-inert-reduction
\chapters@set{20261125T008000}{13}% one-cm-determinant
\chapters@set{20261124T214500}{14}% supersingular-theta-divisibility
\chapters@set{20261124T214600}{15}% the-isotropic-plane
\chapters@set{20261125T013700}{16}% pairings-on-a-formal-group-of-height-two
\chapters@set{20261125T013500}{17}% hecke-nondegeneracy-on-the-tree
\chapters@set{20261126T030200}{18}% the-bound-is-attained
\chapters@set{20261125T009000}{19}% kronecker-limit-and-chowla-selberg
\chapters@set{20261125T002000}{20}% definite-theta-lifts
\chapters@set{20261125T015500}{21}% siegel-weil-in-the-smallest-case
\chapters@set{20261125T002100}{22}% definite-theta-lifts-and-norms
\chapters@set{20261125T016500}{23}% hecke-operators-and-theta-lifts
\chapters@set{20261125T016700}{24}% divisibility-along-the-chain
\chapters@set{20261125T004000}{25}% denominators-and-sha
\chapters@set{20261125T004100}{26}% tate-twists-and-sha
\chapters@set{20261126T020430}{27}% the-minus-part-by-gauss-sums
\chapters@set{20261125T014100}{28}% congruence-modules
\chapters@set{20261125T014500}{29}% eisenstein-congruences-and-selmer-classes
\chapters@set{20261125T014700}{30}% pseudocharacters-and-the-length-bound
\chapters@set{20261125T011000}{31}% galois-representations-black-box
\chapters@set{20261126T030300}{32}% the-data
\chapters@set{20261126T020210}{33}% bessel-poincare-series-and-their-gram-operators
\chapters@set{20261126T030400}{34}% the-seesaw-identity
\chapters@set{20261126T030500}{35}% divisibility-of-the-composite-lift
\chapters@set{20261126T030600}{36}% theorem-a-short-forms
\chapters@set{20261126T030700}{37}% the-determinant-of-a-chain
\chapters@set{20261126T030800}{38}% the-cost-of-one-lift
\chapters@set{20261126T030900}{39}% denominators-of-projections-to-lifts
\chapters@set{20261126T031000}{40}% twists-with-few-congruences
\chapters@set{20261126T031100}{41}% theorem-c-and-the-main-inequality
\providecommand\chapnumber[1]{%
  \@ifundefined{chapters@#1}%
    {\PackageError{chapters}{Sheet #1 is not listed in the index}%
      {List it in 20261125T012000--grh-workbook__guide.org\MessageBreak
       and run scripts/make-chapters.py.}}%
    {\edef\thechap{\@nameuse{chapters@#1}}}}
\makeatother
\chapnumber{<sheet ID>}
\makeatletter
\providecommand\chapters@set[2]{\@namedef{chapters@#1}{#2}}
\chapters@set{20261125T001000}{1}% zeros-to-ramp-sums
\chapters@set{20261125T001100}{2}% inert-primes-and-estimate-b
\chapters@set{20261125T005000}{3}% galois-fourier-and-averaging
\chapters@set{20261125T005500}{4}% columns-from-a-sphere
\chapters@set{20261125T000000}{5}% slope-method
\chapters@set{20261126T030100}{6}% a-twist-for-each-prime
\chapters@set{20261125T003000}{7}% the-chain-and-the-square-root
\chapters@set{20261125T010000}{8}% rank-lemma-and-zero-estimate
\chapters@set{20261125T010500}{9}% poincare-series-and-near-orthogonality
\chapters@set{20261124T223000}{10}% elliptic-curves-for-the-proof
\chapters@set{20261125T009500}{11}% fock-model-and-theta-intertwiners
\chapters@set{20261124T223100}{12}% complex-multiplication-and-inert-reduction
\chapters@set{20261125T008000}{13}% one-cm-determinant
\chapters@set{20261124T214500}{14}% supersingular-theta-divisibility
\chapters@set{20261124T214600}{15}% the-isotropic-plane
\chapters@set{20261125T013700}{16}% pairings-on-a-formal-group-of-height-two
\chapters@set{20261125T013500}{17}% hecke-nondegeneracy-on-the-tree
\chapters@set{20261126T030200}{18}% the-bound-is-attained
\chapters@set{20261125T009000}{19}% kronecker-limit-and-chowla-selberg
\chapters@set{20261125T002000}{20}% definite-theta-lifts
\chapters@set{20261125T015500}{21}% siegel-weil-in-the-smallest-case
\chapters@set{20261125T002100}{22}% definite-theta-lifts-and-norms
\chapters@set{20261125T016500}{23}% hecke-operators-and-theta-lifts
\chapters@set{20261125T016700}{24}% divisibility-along-the-chain
\chapters@set{20261125T004000}{25}% denominators-and-sha
\chapters@set{20261125T004100}{26}% tate-twists-and-sha
\chapters@set{20261126T020430}{27}% the-minus-part-by-gauss-sums
\chapters@set{20261125T014100}{28}% congruence-modules
\chapters@set{20261125T014500}{29}% eisenstein-congruences-and-selmer-classes
\chapters@set{20261125T014700}{30}% pseudocharacters-and-the-length-bound
\chapters@set{20261125T011000}{31}% galois-representations-black-box
\chapters@set{20261126T030300}{32}% the-data
\chapters@set{20261126T020210}{33}% bessel-poincare-series-and-their-gram-operators
\chapters@set{20261126T030400}{34}% the-seesaw-identity
\chapters@set{20261126T030500}{35}% divisibility-of-the-composite-lift
\chapters@set{20261126T030600}{36}% theorem-a-short-forms
\chapters@set{20261126T030700}{37}% the-determinant-of-a-chain
\chapters@set{20261126T030800}{38}% the-cost-of-one-lift
\chapters@set{20261126T030900}{39}% denominators-of-projections-to-lifts
\chapters@set{20261126T031000}{40}% twists-with-few-congruences
\chapters@set{20261126T031100}{41}% theorem-c-and-the-main-inequality
\providecommand\chapnumber[1]{%
  \@ifundefined{chapters@#1}%
    {\PackageError{chapters}{Sheet #1 is not listed in the index}%
      {List it in 20261125T012000--grh-workbook__guide.org\MessageBreak
       and run scripts/make-chapters.py.}}%
    {\edef\thechap{\@nameuse{chapters@#1}}}}
\makeatother
\chapnumber{<sheet ID>}
\makeatletter
\providecommand\chapters@set[2]{\@namedef{chapters@#1}{#2}}
\chapters@set{20261125T001000}{1}% zeros-to-ramp-sums
\chapters@set{20261125T001100}{2}% inert-primes-and-estimate-b
\chapters@set{20261125T005000}{3}% galois-fourier-and-averaging
\chapters@set{20261125T005500}{4}% columns-from-a-sphere
\chapters@set{20261125T000000}{5}% slope-method
\chapters@set{20261126T030100}{6}% a-twist-for-each-prime
\chapters@set{20261125T003000}{7}% the-chain-and-the-square-root
\chapters@set{20261125T010000}{8}% rank-lemma-and-zero-estimate
\chapters@set{20261125T010500}{9}% poincare-series-and-near-orthogonality
\chapters@set{20261124T223000}{10}% elliptic-curves-for-the-proof
\chapters@set{20261125T009500}{11}% fock-model-and-theta-intertwiners
\chapters@set{20261124T223100}{12}% complex-multiplication-and-inert-reduction
\chapters@set{20261125T008000}{13}% one-cm-determinant
\chapters@set{20261124T214500}{14}% supersingular-theta-divisibility
\chapters@set{20261124T214600}{15}% the-isotropic-plane
\chapters@set{20261125T013700}{16}% pairings-on-a-formal-group-of-height-two
\chapters@set{20261125T013500}{17}% hecke-nondegeneracy-on-the-tree
\chapters@set{20261126T030200}{18}% the-bound-is-attained
\chapters@set{20261125T009000}{19}% kronecker-limit-and-chowla-selberg
\chapters@set{20261125T002000}{20}% definite-theta-lifts
\chapters@set{20261125T015500}{21}% siegel-weil-in-the-smallest-case
\chapters@set{20261125T002100}{22}% definite-theta-lifts-and-norms
\chapters@set{20261125T016500}{23}% hecke-operators-and-theta-lifts
\chapters@set{20261125T016700}{24}% divisibility-along-the-chain
\chapters@set{20261125T004000}{25}% denominators-and-sha
\chapters@set{20261125T004100}{26}% tate-twists-and-sha
\chapters@set{20261126T020430}{27}% the-minus-part-by-gauss-sums
\chapters@set{20261125T014100}{28}% congruence-modules
\chapters@set{20261125T014500}{29}% eisenstein-congruences-and-selmer-classes
\chapters@set{20261125T014700}{30}% pseudocharacters-and-the-length-bound
\chapters@set{20261125T011000}{31}% galois-representations-black-box
\chapters@set{20261126T030300}{32}% the-data
\chapters@set{20261126T020210}{33}% bessel-poincare-series-and-their-gram-operators
\chapters@set{20261126T030400}{34}% the-seesaw-identity
\chapters@set{20261126T030500}{35}% divisibility-of-the-composite-lift
\chapters@set{20261126T030600}{36}% theorem-a-short-forms
\chapters@set{20261126T030700}{37}% the-determinant-of-a-chain
\chapters@set{20261126T030800}{38}% the-cost-of-one-lift
\chapters@set{20261126T030900}{39}% denominators-of-projections-to-lifts
\chapters@set{20261126T031000}{40}% twists-with-few-congruences
\chapters@set{20261126T031100}{41}% theorem-c-and-the-main-inequality
\providecommand\chapnumber[1]{%
  \@ifundefined{chapters@#1}%
    {\PackageError{chapters}{Sheet #1 is not listed in the index}%
      {List it in 20261125T012000--grh-workbook__guide.org\MessageBreak
       and run scripts/make-chapters.py.}}%
    {\edef\thechap{\@nameuse{chapters@#1}}}}
\makeatother

\ifcsname chapters@20261126T020110\endcsname
  \chapnumber{20261126T020110}
\else
  \renewcommand{\thechap}{?}
\fi
\renewcommand{\thesection}{\thechap.\arabic{section}}
\theoremstyle{definition}
\newtheorem{problem}{Problem}[section]
\theoremstyle{plain}
\newtheorem{theorem}{Theorem}[section]
\newtheorem{lemma}[theorem]{Lemma}
\theoremstyle{remark}
\newtheorem*{leanremark}{In the Lean formalization}
\newcommand{\Z}{\mathbb Z}
\newcommand{\Q}{\mathbb Q}
\newcommand{\C}{\mathbb C}
\newcommand{\Oc}{\mathcal O}
\newcommand{\Kc}{\mathcal K}
\newcommand{\Ec}{\mathcal E}
\renewcommand{\lg}{\operatorname{lg}}
\newcommand{\ii}{\mathrm i}
\newcommand{\tl}[1]{\texttt{#1}}
\newcommand{\pe}{+_{\Ec}}
\setlength{\emergencystretch}{3em}

\title{The normalized kernel law}
\author{}
\date{November 26, 2026}

\begin{document}
\maketitle
\refstepcounter{chap}\label{chap}

\paragraph{Prerequisites.}
Chapter~\ref{el-chap} on the Weierstrass functions, the quasi-period map \(\eta\) and Legendre's relation; Chapter~\ref{fk-chap} on the operators \(U_x\) and theta functions of a lattice; Chapter~\ref{cm-chap} on formal groups and Lubin--Tate theory; Chapter~\ref{di-chap}, whose notation we keep: \(K=\Q(\sqrt{-d})\) with \(d\equiv7\bmod16\) prime, \(p\ge5\) inert in \(K\), \(W=W(\overline{\mathbb F}_p)\) with Frobenius \(\sigma\), \(\Kc_0=W[1/p]\), \(\lg(X)=\sum_{k\ge0}X^{p^{2k}}/p^k\), and the isomorphism \(\iota\colon\C\to\C_p\).  A series over \(\C_p\) is \emph{integral} if its coefficients have absolute value at most \(1\).  Formal groups and reduction of elliptic curves are as in \cite[Chs.~IV and~VII]{AEC}.  Since \(\sigma\) is Frobenius, we write \(\sigma_\Lambda\) for the Weierstrass sigma function.

Theorem~\ref{di-thm:kernel} writes a theta kernel at an inert prime as \(c\exp(-\sum_e\beta(x_e,y_e))\,U\) with \(U\) integral, and says nothing about the scalar \(c\in\C_p^\times\).  Condition (a) of Theorem~\ref{di-thm:true-lifts}, Adriazola's hypothesis (N3) \cite[Hyp.~11.42]{Adr}, asks that \(\lvert c\rvert\le1\) for the kernels of the chain.  Which normalization of the theta function gives this, and how is it proved?  In the Lean formalization \cite{Lean} every kernel is built from one theta function of one elliptic curve with complex multiplication by \(\Oc_K\), normalized to take the value \(1\) at the origin.  The scalar is then controlled by a fourth tensor power, in which it appears as a coefficient.

\section{One curve and its reduced theta function}\label{sec:one}
Let \(\Lambda\subset\C\) be a lattice with \(\Oc_K\Lambda\subset\Lambda\), and let \(A=\lvert\operatorname{Im}(\bar\omega_1\omega_2)\rvert\) be its covolume, for a basis \(\omega_1,\omega_2\).  For \(\omega,\lambda\in\Lambda\) the number \(\operatorname{Im}(\bar\omega\lambda)\) is an integer multiple of \(A\).  Let \(\eta\colon\Lambda\to\C\) be the quasi-period map, \(\zeta_\Lambda(z+\omega)=\zeta_\Lambda(z)+\eta(\omega)\).

\begin{problem}\label{pr:s2}
There is a unique \(s_2=s_2(\Lambda)\in\C\) with \(\eta(\omega)=s_2\omega+\frac\pi A\bar\omega\) for all \(\omega\in\Lambda\).  Moreover \(s_2(c\Lambda)=c^{-2}s_2(\Lambda)\) for \(c\in\C^\times\).
\end{problem}

The \emph{reduced theta function} is \(\theta(z)=\sigma_\Lambda(z)e^{-s_2z^2/2}\).  It is odd, \(\theta'(0)=1\), its zeros are the points of \(\Lambda\), and \(\theta_{c\Lambda}(cz)=c\,\theta_\Lambda(z)\).  Put \(H(z,w)=z\bar w/A\) and \((U_xG)(z)=G(z+x)\,e^{-\pi H(z,x)-\frac\pi2H(x,x)}\), as in Chapter~\ref{fk-chap}.

\begin{problem}\label{pr:theta-qp}
For \(\omega\in\Lambda\), \(\theta(z+\omega)=\psi(\omega)\,\theta(z)\exp\big(\frac\pi A\bar\omega(z+\frac\omega2)\big)\), with \(\psi\) as in Problem~\ref{el-pr:sigma-qp}.  Equivalently \(\psi(\omega)\,U_\omega\theta=\theta\) for \(\omega\in\Lambda\).
\end{problem}

Let \(t\in\C\setminus\Lambda\) with \(Nt\in\Lambda\) for some integer \(N\ge1\) prime to \(p\).  The \emph{normalized translate} of \(\theta\) by \(t\) is \(B_t=U_t\theta/(U_t\theta)(0)\).

\begin{problem}\label{pr:Bt}
\((U_t\theta)(0)=\theta(t)e^{-\pi\lvert t\rvert^2/(2A)}\ne0\), and \(B_t(z)=\theta(z+t)\,e^{-\pi\bar tz/A}/\theta(t)\); in particular \(B_t(0)=1\).  The quotient \(k_t=B_t/\theta\) satisfies \(k_t(z+\omega)=e^{2\pi\ii\operatorname{Im}(\bar\omega t)/A}\,k_t(z)\) for \(\omega\in\Lambda\), so \(k_t^N\) is \(\Lambda\)-periodic.
\end{problem}

\paragraph{The model at \(p\).}  By the theory of complex multiplication, \(j(\Lambda)\) is an algebraic integer, and there is an algebraic \(u\) with \(\lvert\iota(u)\rvert=1\) such that \(y^2=4x^3-G_2x-G_3\), with \(G_2=u^4g_2(\Lambda)\) and \(G_3=u^6g_3(\Lambda)\), has good reduction over \(W\) \cite[Prop.~8.4]{Adr}.  (Adriazola cites Serre--Tate for this; the Lean derives it from global Artin reciprocity.)  Let \(\Ec\) be this curve over \(W\), \(T=-2x/y\) its parameter at the origin, \(T_1\pe T_2\) its formal group law, and \(\ell(T)=T+\cdots\) its logarithm.  Then \(z=u\,\ell(T)\) as power series \cite[\S8.2]{Adr}.  Put
\[
 \Upsilon(T)=-\log\frac{\theta(u\,\ell(T))}{u\,T}.
\]

\begin{problem}\label{pr:log-theta}
If \(\wp_\Lambda(z)=z^{-2}+\sum_{k\ge2}e_kz^{2k-2}\), then \(\log(\theta(z)/z)=-\frac{s_2}2z^2-\sum_{k\ge2}\frac{e_k}{(2k-1)2k}z^{2k}\).  Hence
\(\Upsilon(T)=-\log\frac{\ell(T)}T+\frac{u^2s_2}2\ell(T)^2+\sum_{k\ge2}\frac{u^{2k}e_k}{(2k-1)2k}\ell(T)^{2k}\).
\end{problem}

The numbers \(u^2s_2\) and \(u^{2k}e_k\) lie in \(\Kc_0\); for \(s_2\) this comes from an expression of \(s_2\) through values of \(\wp_\Lambda\) at torsion points of order prime to \(p\) \cite[Lemma~8.3(iv) and \S8.2]{Adr}.  So \(\Upsilon\) is a series over \(\Kc_0\) without constant term.  It is not integral.  For series \(B,B'\) without constant term we write \(B\simeq B'\) when \(\exp(B-B')\) is integral.

\begin{problem}\label{pr:simeq}
The relation \(\simeq\) is an equivalence relation.  If \(B_1\simeq B_1'\) and \(B_2\simeq B_2'\), then \(B_1+B_2\simeq B_1'+B_2'\), and \(B_1\circ\phi\simeq B_1'\circ\phi\) for every family \(\phi\) of integral series without constant term.
\end{problem}

\begin{theorem}[One variable]\label{thm:one}
For \(t\) as above, \(e^{\Upsilon(T)}B_t(u\ell(T))\) is an integral series with constant term \(1\).
\end{theorem}
\noindent\emph{Status}: proved in \cite[Prop.~8.8]{Adr}.  By Problem~\ref{pr:Bt}, \(k_t^N\) is a rational function on the curve with divisor \(N[-t]-N[0]\); the point \(-t\) does not reduce to the origin, so near the origin of the special fibre \(k_t^N\) is \(T^{-N}\) times a unit, up to a power of a uniformizer.  The normalization \(B_t(0)=1\) removes that power, and an \(N\)-th root is taken by the binomial series, since \(N\) is prime to \(p\).

\section{Many variables on one curve}\label{sec:rules}
For \(\alpha\in\Z_{p^2}\) let \([\alpha]\) be the endomorphism of the formal group of \(\Ec\) with \(\ell([\alpha]T)=\alpha\,\ell(T)\); it has coefficients in \(W\) by Lubin--Tate theory, as in Chapter~\ref{cm-chap}.  On \(\Z_{p^2}\) the Frobenius \(\sigma\) is the nontrivial automorphism, and \(\sigma(\iota(a))=\iota(\bar a)\) for \(a\in\Oc_K\).  Put \(\delta\Upsilon(S,T)=\Upsilon(S\pe T)-\Upsilon(S)-\Upsilon(T)\).

\begin{theorem}[Biextension rules]\label{thm:rules}
For all \(\alpha\in\Z_{p^2}\):
\begin{enumerate}
\item[(R1)] \(\delta\Upsilon(S\pe S',T)\simeq\delta\Upsilon(S,T)+\delta\Upsilon(S',T)\), and the same in the second variable;
\item[(R2)] \(\delta\Upsilon([\alpha]S,T)\simeq\delta\Upsilon(S,[\sigma\alpha]T)\);
\item[(R3)] \(\Upsilon([\alpha]T)\simeq\alpha\,\sigma(\alpha)\,\Upsilon(T)\).
\end{enumerate}
\end{theorem}
\noindent\emph{Status}: proved in \cite[Props.~8.9, 8.10 and Lemma~8.13]{Adr}.  Rule (R1) is the theorem of the cube: the quotient \(\theta(z_1+z_2+z_3)\theta(z_1)\theta(z_2)\theta(z_3)/\theta(z_1+z_2)\theta(z_1+z_3)\theta(z_2+z_3)\) is a rational function on the cube of the curve whose divisor near the origin consists of kernels of homomorphisms with a unit coefficient (compare Problem~\ref{el-pr:cube}).  Rules (R2) and (R3) come in the same way from quotients of \(\theta\) at \(az\) and \(\bar aw\) for \(a\in\Oc_K\), and extend to \(\Z_{p^2}\) by continuity.  Only complex multiplication by \(\Oc_K\) enters.

\begin{problem}\label{pr:diag}
For every \(c\in\Z_{p^2}\), \(\delta\Upsilon(T,[c]T)\simeq(c+\sigma(c))\,\Upsilon(T)\).
\end{problem}

\begin{problem}\label{pr:upssum}
For integral series \(A_1,\dots,A_k\) without constant term,
\(\Upsilon(A_1\pe\cdots\pe A_k)\simeq\sum_j\Upsilon(A_j)+\sum_{j<l}\delta\Upsilon(A_j,A_l)\).
\end{problem}

\begin{problem}\label{pr:unitary}
Let \(\gamma\in M_m(\Z_{p^2})\) satisfy \(\sum_i\sigma(\gamma_{ij})\gamma_{il}=[j=l]\), and put \(T'_i=[\gamma_{i1}]T_1\pe\cdots\pe[\gamma_{im}]T_m\).  Then \(\sum_i\Upsilon(T'_i)\simeq\sum_j\Upsilon(T_j)\).
\end{problem}

A \emph{hyperbolic pair} of size \(m\) is a pair of \(m\times2m\) matrices \(f,f'\) over \(\Z_{p^2}\) with
\(\sum_i\sigma(f_{e,i})f_{e',i}=0\), \(\sum_i\sigma(f'_{e,i})f'_{e',i}=0\) and \(\sum_i\sigma(f_{e,i})f'_{e',i}=[e=e']\).

\begin{problem}\label{pr:hypshape}
For a hyperbolic pair of size \(m\) put \(T_i=[f_{1,i}]X_1\pe\cdots\pe[f_{m,i}]X_m\pe[f'_{1,i}]Y_1\pe\cdots\pe[f'_{m,i}]Y_m\) for \(i\le2m\).  Then \(\sum_{i\le2m}\Upsilon(T_i)\simeq\sum_{e\le m}\delta\Upsilon(X_e,Y_e)\).
\end{problem}

Let \(h\) be a hermitian form on a \(\Z_{p^2}\)-module, linear in the first variable, so \(h(y,x)=\sigma(h(x,y))\).  Since \(\Q_{p^2}/\Q_p\) is unramified, the norm \(\kappa\mapsto\kappa\sigma(\kappa)\) maps \(\Z_{p^2}^\times\) onto \(\Z_p^\times\).

\begin{problem}\label{pr:hypplane}
Let \(h(f,f)=h(f',f')=0\), \(h(f,f')=1\), and \(\kappa\sigma(\kappa)=-1\).  Then \(e_1=f+\frac12f'\) and \(e_2=\kappa(f-\frac12f')\) satisfy \(h(e_i,e_j)=[i=j]\) and span the same \(\Z_{p^2}\)-module as \(f,f'\).
\end{problem}

\paragraph{The coboundary in Lubin--Tate coordinates.}  The formal group of \(\Ec\) and the group with logarithm \(\lg\) are isomorphic over \(W\): there are \(\varepsilon\in W^\times\) and \(\varkappa\in W[[T]]\) with \(\varkappa'(0)\in W^\times\) and \(\lg(\varkappa(T))=\varepsilon\,\ell(T)\) \cite[Prop.~8.11]{Adr}, \cite{LT}.  Put \(\beta_0(x,y)=\delta\Upsilon(\varkappa^{-1}x,\varkappa^{-1}y)\).  This is the series \(\beta\) of Chapter~\ref{di-chap}.

The Lean uses an averaged form of \(\beta_0\).  Let \(2\nu\) be the length of a Frobenius orbit, as in Theorem~\ref{di-thm:class}, with one variable \(X_s\) for each position \(s\in\Z/2\nu\).  The torus \(\mu_{p^{2\nu}-1}\) acts by \(X_s\mapsto \zeta^{p^s}X_s\) for even \(s\) and \(X_s\mapsto \zeta^{-p^s}X_s\) for odd \(s\), where \(0\le s<2\nu\).

\begin{problem}\label{pr:torus}
\begin{enumerate}
\item[(i)] For \(0\le s,t<2\nu\), \(p^s\equiv p^t\pmod{p^{2\nu}-1}\) only if \(s=t\).
\item[(ii)] The average of a series over the torus keeps exactly the monomials \(\prod_sX_s^{m_s}\) with \(\sum_{s\ \mathrm{even}}m_sp^s-\sum_{s\ \mathrm{odd}}m_sp^s\equiv0\pmod{p^{2\nu}-1}\).
\item[(iii)] It removes every monomial \(X_sX_t\) with \(s\) even and \(t\) odd, and keeps \(X_s^pX_{s+1}\) for every \(s\in\Z/2\nu\).
\end{enumerate}
\end{problem}

The monomials \(X_s^pX_{s+1}\) are those of the class \(h_\upsilon\) of Section~\ref{di-sec:minors}.  In the Lean, \(\beta\) is the torus average of \(\sum_r\beta_0\big(x_r,y_r\big)\), where \(x_r,y_r\) are the Lubin--Tate coordinates obtained from the trace-dual bases \(\omega_r,\omega_r^*\) as in Theorem~\ref{di-thm:class}.  It satisfies four conditions: (B1) its coefficients lie in \(\Kc_0\); (B2) it is invariant under the torus; (B3) each of its monomials has positive degree in the even and in the odd variables and is not of bidegree \((1,1)\); (B4) the exponentials of its coboundaries in the even and in the odd variables, taken in logarithmic coordinates, are integral with constant term \(1\).  \emph{Status}: proved in the Lean formalization from (R1)--(R3) and Hazewinkel's functional-equation lemma.  The Lean also proves that replacing the exponent by its torus average changes the kernel only by an integral factor; we have not followed that argument.  Adriazola works with \(\beta_0\) itself \cite[Thm.~8.12]{Adr}.

\section{The scalar}\label{sec:scalar}
Write \(f^{\otimes4}\) for the product \(f(x^{(1)})f(x^{(2)})f(x^{(3)})f(x^{(4)})\) of four copies of a series \(f(x_1,\dots,x_n)\) in disjoint sets of variables.

\begin{problem}\label{pr:fourth}
If every coefficient of \(f^{\otimes4}\) has absolute value at most \(R\), then every coefficient of \(f\) has absolute value at most \(R^{1/4}\).  In particular, if \(f^{\otimes4}=\kappa g\) with \(g\) integral, \(g(0)=1\) and \(\lvert\kappa\rvert\le1\), then \(f\) is integral and \(\lvert f(0)\rvert=\lvert\kappa\rvert^{1/4}\).
\end{problem}

Adriazola uses the same idea with two copies.  For a theta function \(G\) of a hermitian lattice he writes \(G\otimes G\), after a change of coordinates over \(K\), as a combination \(\sum_tr_tB_t\) of products of normalized translates, and obtains a bound \(R=\max_t\lvert r_t\rvert\) for the coefficients of the tensor square \cite[Thm.~8.19]{Adr}.  Nothing controls the \(r_t\), so the conclusion is \(c\,U\) with an unknown scalar \(c\).  The Lean fixes the normalization of the theta function so that the corresponding coefficients are integral.

By Landherr's theorem, positive definite hermitian spaces over \(K\) of equal dimension whose discriminants differ by a norm from \(K\) are isometric \cite[Ch.~10, \S6]{Sch}.

\begin{problem}\label{pr:landherr}
The discriminant of the sum of \(k\) copies of a hermitian space over \(K\) is the \(k\)-th power of its discriminant, and for even \(k\) it is a norm from \(K\).  Hence the sum of four copies of a positive definite hermitian space over \(K\) is isometric to the standard space \(\sum_ix_i\bar y_i\) of the same dimension.
\end{problem}

\paragraph{The kernels.}  Let \(E\supset K\) be the CM field of Chapter~\ref{di-chap}, \(\Phi\) its CM type, and \(\mathfrak a\) the \(\Oc_E\)-ideal used there for the hermitian lattices.  Take \(\Lambda=\delta\Oc_K\) with \(\delta^{12}=g_2(\Oc_K)^3-27g_3(\Oc_K)^2\), so that \(g_2^3-27g_3^2=1\) for \(\Lambda\).  For even \(N,N'\) and self-dual \(\Oc_E\)-lattices \(L\subset E^N\), \(L'\subset E^{N'}\), let \(Y(L,L')=\mathfrak a\operatorname{Hom}(L,L')\) with the \(K\)-valued form \(\operatorname{Tr}_{E/K}\sum_{i,j}x_{ij}\bar y_{ij}\).  Its theta functions live on \(Y(L,L')\otimes_{\Oc_K}\C\), with the lattice \(Y(L,L')\otimes_{\Oc_K}\Lambda\); only the curve \(\C/\Lambda\) enters.  The construction has three layers.
\begin{enumerate}
\item[(a)] A theta function \(\theta_{\mathfrak a}\) on \(\mathfrak a\otimes\C=\C^\Phi\) for the lattice \(\mathfrak a\otimes\Lambda\).  Fix a \(K\)-linear isometry \(\gamma\) from four copies of \(\mathfrak a\) to the standard hermitian space of the same dimension (Problem~\ref{pr:landherr}).  The function \(\theta_{\mathfrak a}\) is defined by
\[
 \theta_{\mathfrak a}^{\otimes4}=\kappa\cdot I_\gamma\big(B_{t_0}^{\otimes4\lvert\Phi\rvert}\big)\circ\gamma ,
\]
up to a fourth root of unity.  Here \(B_{t_0}\) is the normalized translate by a point \(t_0\) of order two, \(I_\gamma\) is the average of the translates \(U_x\) over the cosets of the standard lattice in the image under \(\gamma\) of four copies of \(\mathfrak a\otimes\Lambda\), weighted by a semicharacter and divided by the square root of the index, and \(\kappa\) is an algebraic number.  For each inert \(p\) the same function satisfies \(\theta_{\mathfrak a}^{\otimes4}=\kappa_p\cdot I_{\gamma_p}(B_{t_0}^{\otimes4\lvert\Phi\rvert})\circ\gamma_p\) for an isometry \(\gamma_p\) whose index is prime to \(p\), with \(\lvert\iota(\kappa_p)\rvert\le1\).
\item[(b)] On standard lattices, the product of copies of \(\theta_{\mathfrak a}\), one for each matrix entry.
\item[(c)] For general \((L,L')\), coset averages of (b), with index factors and a character factor.
\end{enumerate}
The \emph{kernels} at \((L,L')\) are the combinations \(\sum_kr_kU_{t_k}e_k\) with \(\lvert r_k\rvert\le1\), \(t_k\) torsion of order prime to \(p\), and \(e_k\) the function of (c) for a pair equal to \((L,L')\) at \(p\).

\begin{theorem}[Kernel law in orthonormal coordinates]\label{thm:A}
Let \(b_1,\dots,b_r\) be an orthonormal basis of \(Y(L,L')\otimes_{\Oc_K}\Z_{p^2}\).  For every kernel \(F\) at \((L,L')\),
\[
 \exp\Big(\sum_{i=1}^r\Upsilon(T_i)\Big)\cdot F\Big(\sum_{i=1}^rb_i\otimes u\,\ell(T_i)\Big)
\]
is integral.  There is no scalar.
\end{theorem}
\noindent As in Chapter~\ref{di-chap}, \(F\) stands for its Taylor series moved to \(\C_p\) by \(\iota\), and \(b_i\) is substituted through \(\Oc_K\subset\Z_{p^2}\subset W\).
\noindent\emph{Status}: proved in the Lean formalization.  The proof is the following chain; each step is proved there.
\begin{enumerate}
\item[(1)] \emph{One copy.}  For \(\theta_{\mathfrak a}\) and its translates by torsion of order prime to \(p\), the assertion holds in one orthonormal frame of \(\mathfrak a\otimes\Z_{p^2}\).  By (a), after a change of variables by a \(\sigma\)-unitary matrix over \(W\), the fourth tensor power becomes \(\kappa_p\) times a combination \(\sum_kr_k\prod_wB_{t_{k,w}}\) with \(\lvert r_k\rvert\le1\).  Theorem~\ref{thm:one} handles each factor, Problem~\ref{pr:unitary} (applied in four copies) handles the change of variables, and Problem~\ref{pr:fourth} takes the fourth root.  That the values at \(0\) entering the \(r_k\) are \(p\)-integral is proved in the Lean; we have not followed that proof.
\item[(2)] \emph{Products.}  The case of (b) follows, since integrality is preserved by products in disjoint variables.
\item[(3)] \emph{Change of frame.}  The assertion for one orthonormal frame implies it for every orthonormal frame (Problem~\ref{pr:frame}).
\item[(4)] \emph{Standard at \(p\).}  If \((M,M')\) equals the standard pair at \(p\), every \(U_te(M,M')\) is a \(p\)-integral combination of translates of (b).  We have not followed this step, which is a computation with the coset averages (c).
\item[(5)] \emph{Weak approximation and covariance.}  Every pair is \((\gamma M,\gamma'M')\) with \(\gamma,\gamma'\) unitary over \(E\) and \((M,M')\) standard at \(p\), and the kernel of \((\gamma M,\gamma'M')\) is the kernel of \((M,M')\) composed with \((\gamma,\gamma')\).
\item[(6)] \emph{Combinations.}  The assertion is stable under combinations with integral coefficients and depends on \((L,L')\) only at \(p\).
\end{enumerate}

\begin{problem}\label{pr:frame}
Let \(b'_j=\sum_i\gamma_{ij}b_i\) with both bases orthonormal.  Then \(\sum_i\sigma(\gamma_{ij})\gamma_{il}=[j=l]\), and \(\sum_jb'_j\otimes u\ell(T_j)=\sum_ib_i\otimes u\ell(T'_i)\) with \(T'_i\) as in Problem~\ref{pr:unitary}.  Deduce that Theorem~\ref{thm:A} for the basis \(b\) implies it for \(b'\).
\end{problem}

\begin{theorem}[Normalized kernel law, (N3)]\label{thm:B}
Let \(F\) be a kernel at \((L,L')\), and suppose \(Y(L,L')\otimes\Z_{p^2}=P\oplus P^\perp\) with \(P^\perp\) orthonormal and \(P\) spanned by \(f_1,\dots,f_m,f'_1,\dots,f'_m\) with \(h(f_e,f_{e'})=h(f'_e,f'_{e'})=0\) and \(h(f_e,f'_{e'})=[e=e']\).  Put \(\Theta(x,y)=F\big(\sum_ef_e\otimes u\varepsilon^{-1}\lg(x_e)+\sum_ef'_e\otimes u\varepsilon^{-1}\lg(y_e)\big)\), in the coordinates of Theorem~\ref{di-thm:kernel}.  Then
\[
 \Theta(x,y)=\exp\Big(-\sum_{e=1}^m\beta_0(x_e,y_e)\Big)\,U(x,y),\qquad U\ \text{integral}.
\]
\end{theorem}

\begin{problem}\label{pr:N3}
Theorem~\ref{thm:A}, Problems~\ref{pr:hypplane}, \ref{pr:hypshape} and~\ref{pr:simeq}, and the relation \(\lg\circ\varkappa=\varepsilon\ell\) imply Theorem~\ref{thm:B}.  Thus the scalar of Theorem~\ref{di-thm:kernel} can be taken to be \(1\) for the kernels of the chain, and Theorem~\ref{di-thm:class} holds with \(c=1\).
\end{problem}

Theorem~\ref{thm:B} is Theorem~\ref{di-thm:kernel} with \(c=1\), for the kernels just described, which are the kernels whose jets define the lifts of the chain.  The step kernels of Chapter~\ref{di-chap} carry in addition the scalar \(s_j\) of the step normalization, which is a unit; this is (N1), treated in Chapter~\ref{fo-chap}.  \emph{Status}: proved in the Lean formalization.  The Lean proves Theorem~\ref{thm:A} at every admissible prime \(p\ge5\) inert in \(K\), and states the hyperbolic form in the coordinates of the places at the primes whose Frobenius orbits have even length at least \(2\).  Where Frobenius acts on the CM type as complex conjugation, the orbits have length \(1\), there is no hyperbolic splitting along an orbit, and the Lean uses a weaker shape that gives no gain.

\begin{leanremark}\sloppy
The one-variable function is \tl{Formalism.L2.Th1.F.th0} (the translate of \tl{sgStar} divided by its value at \(0\)), with \tl{s2}.  The kernels are the class \tl{IsComb}, built from \tl{Formalism.L2.Ker.Cd.Base.thMM} (characterized by \tl{IsThM}, with \(\kappa\) given by \tl{KapChoice.kappa}), \tl{thZero}, \tl{thMoved} and \tl{eNat}.  Theorem~\ref{thm:A} is \tl{Formalism.L2.KerInert.IntInert}, proved as \tl{kernels\_integral\_inert} through \tl{intInert\_of} from \tl{CombFact}, \tl{PairFact}, \tl{WeakApprox}, \tl{PullFact}, \tl{StdStep}, \tl{TensFact} and \tl{ThFourInt}; the one-copy step is \tl{thFourInt\_of} (from \tl{OneCopyLoc}, \tl{FourProd}, \tl{LocInv}, \tl{GlobONB}), with \tl{scale\_integrality}, \tl{Formalism.L2.Th1.Top.kappa\_integral}, \tl{ThetaFourthTensorPowerProductShape} and \tl{Formalism.L2.Th1.E.FourthPower}.  It is repackaged as \tl{GlobalTheory.KerStd.HostAt} by \tl{hostAt\_inert}.  The rules (R1)--(R3) are \tl{first\_level\_curve\_holds} and \tl{NormRel}; (B1)--(B4) are \tl{B1}--\tl{B4} in \tl{product\_zeroth\_first\_parts}, proved through \tl{biextension\_hypotheses\_of\_level\_one}.  Theorem~\ref{thm:B}, in the coordinates of the places and with the averaged \(\beta\), is \tl{OrbitPlaces.EvenConditionsInert} (with \tl{ThetaExpansionAllPlaces} and \tl{BeOK}), proved by \tl{GlobalTheory.KerStd.exists\_beFam\_evenConditionsInert\_of\_kerEnt}.  The good model is \tl{Formalism.L2.Th1.CM.CM2}, proved by \tl{cmGoodModel\_of\_globalArtinReciprocity} \cite{Lean}.
\end{leanremark}

\section*{Hints}
\noindent\emph{Section~\ref{sec:one}.}
\ref{pr:s2}. Extend \(\eta\) \(\mathbb R\)-linearly to \(\C\) and use Problem~\ref{el-pr:legendre}.
\ref{pr:theta-qp}. Problems~\ref{el-pr:sigma-qp} and~\ref{pr:s2}.
\ref{pr:Bt}. Problem~\ref{pr:theta-qp}, and \(\operatorname{Im}(\bar\omega\lambda)\in A\Z\).
\ref{pr:log-theta}. Differentiate \(\log(\sigma_\Lambda(z)/z)\) twice; it is even.
\ref{pr:simeq}. An integral series with constant term \(1\) has an integral inverse.

\noindent\emph{Section~\ref{sec:rules}.}
\ref{pr:diag}. \(T\pe[c]T=[1+c]T\).
\ref{pr:upssum}. Induction on \(k\), with (R1).
\ref{pr:unitary}. Problem~\ref{pr:upssum}; then (R3) for the diagonal terms and (R2), (R1) for the cross terms.
\ref{pr:hypshape}. As in Problem~\ref{pr:unitary}, with the summands ordered \(X_1,\dots,X_m,Y_1,\dots,Y_m\).
\ref{pr:hypplane}. Expand by sesquilinearity; compute a determinant.
\ref{pr:torus}. For (ii), \(\sum_{\zeta\in\mu_n}\zeta^k\) is \(n\) or \(0\).

\noindent\emph{Section~\ref{sec:scalar}.}
\ref{pr:fourth}. Look at the coefficient of \(\prod_{a\le4}(x^{(a)})^\mu\).
\ref{pr:frame}. Compute \(h(b'_j,b'_l)\); then substitute.
\ref{pr:N3}. Write \(f\otimes u\ell(X)+f'\otimes u\ell(Y)\) in the basis \(e_1,e_2\) of Problem~\ref{pr:hypplane}; the coefficients form a hyperbolic pair.  Then Theorem~\ref{thm:A} with the variables of \(P^\perp\) set to \(0\), Problem~\ref{pr:hypshape}, and \(X_e=\varkappa^{-1}(x_e)\), \(Y_e=\varkappa^{-1}(y_e)\).

\begin{thebibliography}{Lean}
\bibitem[Adr]{Adr}
J.~Adriazola, \emph{The Alp\"oge proof of the Riemann hypothesis}, manuscript, revised 24 November 2026, \url{https://github.com/jimjamzola/RH-Alpoge-Reconstruction}.
\bibitem[AEC]{AEC} J.~H.~Silverman, \emph{The Arithmetic of Elliptic Curves}, 2nd ed., Springer, 2009.
\bibitem[Lean]{Lean} The Lean community, \emph{riemann-hypothesis}, a Lean formalization of Alp\"oge's argument, \url{https://github.com/leanprover-community/riemann-hypothesis}.
\bibitem[LT]{LT} J.~Lubin and J.~Tate, Formal complex multiplication in local fields, Ann. of Math. (2) \textbf{81} (1965), 380--387.
\bibitem[Sch]{Sch} W.~Scharlau, \emph{Quadratic and Hermitian Forms}, Springer, 1985.
\end{thebibliography}
\end{document}
